% COLLECTION_METADATA: {"schema_version": 1, "kind": "reused_writeup", "independently_reviewed": false, "primary_source": "attacks/erdos/gpt_pro_5.2/70.tex", "source_paths": ["attacks/erdos/gpt_5.2/70.tex", "attacks/erdos/gpt_pro_5.2/70.tex"], "source_model": "gpt_pro_5.2", "source_completion": null, "source_claim": "unresolved"}
\section*{Erd\H{o}s Problem 70 --- collection record}

\subsection*{PROVENANCE AND STATUS}
Official problem: https://www.erdosproblems.com/70

Saved collaborative-database status: open.
Snapshot checked: 2026-09-23T17:22:24Z.
Latest status and formalization links: https://teorth.github.io/erdosproblems/

This is an attributed collection of existing writeups. The model-folder names below identify their original attribution; placement in GPT\_6\_Astra\_Ultra does not attribute their mathematical work to that model. All locally available source versions selected for this problem are reproduced below, without editing their text.

Proof claims, computations, literature statements, and completion estimates in the reused text have not been independently verified in this collection pass. A source claiming a solution does not change the saved upstream status.

Primary source (latest numbered version in the preferred model folder): \texttt{attacks/erdos/gpt\_pro\_5.2/70.tex}.

\subsection*{REUSED SOURCE 1}
Original attribution: \texttt{gpt\_5.2}.
Original file: \texttt{attacks/erdos/gpt\_5.2/70.tex}.

% BEGIN REUSED SOURCE: attacks/erdos/gpt_5.2/70.tex

OUTPUT
1) FORMAL RESTATEMENT
Definitions/Conventions.
- Write $\mathfrak c:=|{\mathbb R}|$ for the cardinality of the continuum, and identify it with its initial ordinal.
- For a set $X$ and integer $r\ge 1$, write $[X]^r$ for the family of $r$-element subsets of $X$.
- For cardinals/ordinals $\kappa,\alpha,\gamma$ and $r\ge 1$, the partition relation
\[\kappa \to (\alpha,\gamma)^r_2\]
means: for every coloring $c:[\kappa]^r\to\{0,1\}$, there exists either
  (i) $A\subseteq \kappa$ of order type $\alpha$ such that $c$ is constantly $0$ on $[A]^r$, or
  (ii) $B\subseteq \kappa$ of order type $\gamma$ such that $c$ is constantly $1$ on $[B]^r$.
Here, when $\gamma=n\in\mathbb N$, ``order type $n$'' just means ``a subset of size $n$''.

Question.
Fix a countable ordinal $\beta$ and an integer $n$ with $2\le n<\omega$. Is it true that
\[\mathfrak c \to (\beta,n)^3_2?\]

2) QUICK LITERATURE/CONTEXT CHECK
Web browsing is not available in this session.
From the problem text only:
- Erd\H{o}s and Rado proved $\mathfrak c\to(\omega+n,4)^3_2$ for $2\le n<\omega$.
I do not reproduce that proof here.

3) ATTACK PLAN
Proof track ideas:
1. Reduce to the first nontrivial finite case $n=4$ by monotonicity in the second parameter (proved below), then attempt to extend the known $(\omega+n,4)$ case to all countable $\beta$.
2. Try to build an increasing $\omega_1$-sequence inside $\mathfrak c$ via a transfinite recursion and apply canonical partition theorems.

Disproof track ideas:
1. Try to construct a 2-coloring of $[\mathfrak c]^3$ with no 1-homogeneous $n$-set for some fixed $n\ge 4$ and no 0-homogeneous set of order type $\beta$.
2. Identify which additional axioms (CH, MA, large cardinals) might affect $\mathfrak c\to(\beta,n)^3_2$.

Chosen path in this attempt: isolate trivial cases and monotonicity reductions (fully proved), to clarify the first genuinely open parameter regime.

4) WORK
Lemma 1 (Triviality for $n=2$).
For any infinite cardinal $\kappa$ and any ordinal $\beta$ with $\beta\le \kappa$, one has
\[\kappa\to(\beta,2)^3_2.\]

Proof.
Let $c:[\kappa]^3\to\{0,1\}$ be any coloring. Take any two-element subset $B\subseteq\kappa$.
Then $[B]^3=\emptyset$, so $c$ is vacuously constant $1$ on $[B]^3$.
Thus the second alternative in the definition of $\kappa\to(\beta,2)^3_2$ always holds (provided $\kappa$ has at least two elements, which is true for infinite $\kappa$). \qed

Lemma 2 (Triviality for $n=3$).
For any infinite cardinal $\kappa$ and any ordinal $\beta$ with $\beta\le \kappa$, one has
\[\kappa\to(\beta,3)^3_2.\]

Proof.
Let $c:[\kappa]^3\to\{0,1\}$ be any coloring.
If there exists a triple $B\in[\kappa]^3$ with $c(B)=1$, then $B$ itself witnesses the second alternative, since $[B]^3=\{B\}$.
Otherwise, $c$ is identically $0$ on $[\kappa]^3$. In that case, choose any subset $A\subseteq \kappa$ of order type $\beta$ (possible since $\beta\le\kappa$),
and then $c$ is constantly $0$ on $[A]^3$. So one of the two alternatives always holds. \qed

Lemma 3 (Monotonicity in the finite parameter).
Fix $\kappa,\beta$ and integers $3\le m\le n$. If
\[\kappa\to(\beta,n)^3_2,\]
then
\[\kappa\to(\beta,m)^3_2.\]

Proof.
Let $c:[\kappa]^3\to\{0,1\}$ be any coloring.
Apply $\kappa\to(\beta,n)^3_2$ to obtain either:
(i) a $0$-homogeneous set $A$ of order type $\beta$, in which case we are done; or
(ii) a $1$-homogeneous set $B$ of size $n$ (so $c$ is constantly $1$ on $[B]^3$).
In case (ii), choose any subset $B'\subseteq B$ of size $m$.
Then $[B']^3\subseteq[B]^3$, so $c$ is also constantly $1$ on $[B']^3$, giving the second alternative for $(\beta,m)^3_2$.
Thus $\kappa\to(\beta,m)^3_2$ holds. \qed

FAST REALITY CHECK
The above shows the statement is vacuous for $n=2$ and $n=3$.
So the first nontrivial finite case is $n=4$, and the general question for finite $n$ reduces (by Lemma 3) to understanding the hardest values of $n$.

5) VERIFICATION
- Lemmas 1--2 rely only on the fact that $[B]^3$ is empty for $|B|<3$ and a singleton for $|B|=3$.
- Lemma 3 is a direct subset argument: $1$-homogeneity is inherited by subsets.

6) FINAL
**UNRESOLVED**

(i) Strongest fully proved partial result:
The relation $\kappa\to(\beta,n)^3_2$ is trivial for $n\in\{2,3\}$ (Lemmas 1--2), and is monotone in $n$ (Lemma 3).

(ii) First gap:
Handle the first genuinely nontrivial case: for fixed countable $\beta$, prove or refute
\[\mathfrak c\to(\beta,4)^3_2.\]
The known result in the text only covers $\beta=\omega+n$ with $n=4$.

(iii) Top 3 next moves (concrete):
1. Specialize to $\beta=\omega_1$ (or to $\beta=\omega^\omega$) and try to either extend the Erd\H{o}s--Rado method or construct an explicit bad coloring of $[\mathfrak c]^3$.
2. Determine which set-theoretic axioms affect the truth of $\mathfrak c\to(\beta,n)^3_2$; attempt to prove consistency/independence results for some $\beta,n$.
3. Search for a reduction to canonical colorings on $[\mathfrak c]^3$ (e.g. via a definable well-order of the reals) and analyze homogeneous sets of prescribed order types.

(iv) What a minimal counterexample would likely look like (structure/parameters):
A 2-coloring $c:[\mathfrak c]^3\to\{0,1\}$ such that every $1$-homogeneous set has size $\le n-1$ (no blue $n$-set),
and every $0$-homogeneous set has order type strictly less than the target countable $\beta$.


% END REUSED SOURCE: attacks/erdos/gpt_5.2/70.tex

\subsection*{REUSED SOURCE 2}
Original attribution: \texttt{gpt\_pro\_5.2}.
Original file: \texttt{attacks/erdos/gpt\_pro\_5.2/70.tex}.

% BEGIN REUSED SOURCE: attacks/erdos/gpt_pro_5.2/70.tex

\noindent\textbf{FORMAL RESTATEMENT.}
Let $\mathfrak c$ denote the initial ordinal of cardinality $2^{\aleph_0}$ (the cardinality of $\mathbb R$). Let $\beta$ be a countable ordinal and let $n$ be an integer with $2\le n<\omega$.

For an ordinal $\kappa$ and integer $r\ge 1$, write $[\kappa]^r$ for the set of $r$-element subsets of $\kappa$.
For $\kappa,\beta$ ordinals and $n\in\mathbb N$, the partition relation
\[\kappa\to(\beta,n)_2^3\]
means:

\emph{for every} coloring $c:[\kappa]^3\to\{0,1\}$, \emph{either}
\begin{itemize}
\item there exists $H\subseteq\kappa$ of order type $\beta$ such that $c$ is constantly $0$ on $[H]^3$, \emph{or}
\item there exists $K\subseteq\kappa$ with $|K|=n$ such that $c$ is constantly $1$ on $[K]^3$.
\end{itemize}

Question: is it true that $\mathfrak c\to(\beta,n)_2^3$ for every countable ordinal $\beta$ and every integer $2\le n<\omega$?

\medskip
\noindent\textbf{QUICK LITERATURE/CONTEXT CHECK.}
The problem statement records that Erd\H{o}s and Rado proved
\[\mathfrak c\to(\omega+n,4)_2^3\qquad\text{for every integer }2\le n<\omega.
\]
No additional literature claims are used here.

\medskip
\noindent\textbf{ATTACK PLAN.}
\emph{Proof track ideas.}
(1) Use monotonicity and the known Erd\H{o}s--Rado instances as base cases; attempt to extend from $\beta=\omega+n$ to general countable $\beta$ by transfinite induction on $\beta$.
(2) Try to rework the Erd\H{o}s--Rado proof to handle larger countable order types by a refined fusion/diagonalization.

\emph{Disproof track ideas.}
(1) Attempt to build an explicit $2$-coloring of $[\mathfrak c]^3$ that destroys all $0$-homogeneous sets of a given countable order type while also avoiding $1$-homogeneous $n$-sets.
(2) Check whether the statement might depend on additional axioms (e.g. CH vs. not-CH) by testing known ``pathological'' colorings.

I did not obtain a full proof or counterexample. The WORK section records basic (but problem-specific) reductions.

\medskip
\noindent\textbf{WORK.}

\medskip
\noindent\textbf{Lemma 70.1 (easy cases $n=2$ and $n=3$).}
Let $\kappa$ be an ordinal and $\beta\le\kappa$ an ordinal.
\begin{itemize}
\item[(a)] $\kappa\to(\beta,2)_2^3$ always holds.
\item[(b)] $\kappa\to(\beta,3)_2^3$ always holds.
\end{itemize}
In particular, since every countable ordinal $\beta$ satisfies $\beta<\mathfrak c$, the desired relation holds for all countable $\beta$ when $n\in\{2,3\}$.

\noindent\textbf{Proof.}
(a) Fix any coloring $c:[\kappa]^3\to\{0,1\}$. Choose any $2$-element subset $K\subseteq\kappa$. Then $[K]^3=\varnothing$, so $c$ is vacuously constantly $1$ on $[K]^3$. Thus the second alternative in the definition of $\kappa\to(\beta,2)_2^3$ always holds.

(b) Fix any coloring $c:[\kappa]^3\to\{0,1\}$. If there exists a triple $\{x,y,z\}\in[\kappa]^3$ with $c(\{x,y,z\})=1$, then letting $K=\{x,y,z\}$ gives $|K|=3$ and $c$ is constantly $1$ on $[K]^3$ (since $[K]^3$ contains exactly that one triple). Otherwise, $c$ colors every triple in $[\kappa]^3$ by $0$. In this case, since $\beta\le\kappa$, there exists $H\subseteq\kappa$ of order type $\beta$ (for example, take the first $\beta$ ordinals below $\kappa$), and then $c$ is constantly $0$ on $[H]^3$ by assumption. Hence one of the two alternatives always holds, proving $\kappa\to(\beta,3)_2^3$.
\hfill$\square$

\medskip
\noindent\textbf{Lemma 70.2 (monotonicity in parameters).}
Let $\kappa$ be an ordinal.
\begin{itemize}
\item[(a)] If $\kappa\to(\beta,n)_2^3$ holds and $\beta'\le\beta$, $n'\le n$, then $\kappa\to(\beta',n')_2^3$ holds.
\item[(b)] If $\kappa\to(\beta,n)_2^3$ holds and $\kappa'\ge \kappa$ is an ordinal, then $\kappa'\to(\beta,n)_2^3$ holds.
\end{itemize}

\noindent\textbf{Proof.}
(a) Fix a coloring $c:[\kappa]^3\to\{0,1\}$. By $\kappa\to(\beta,n)_2^3$, either there is a $0$-homogeneous set $H$ of order type $\beta$, or a $1$-homogeneous set $K$ of size $n$.

If the first alternative occurs, then $H$ contains a subset $H'\subseteq H$ of order type $\beta'$ (because $\beta'\le\beta$ and order types are inherited by appropriate subsets). Since $c$ is constantly $0$ on $[H]^3$, it is also constantly $0$ on $[H']^3$.

If the second alternative occurs, choose any subset $K'\subseteq K$ with $|K'|=n'$. Since $[K']^3\subseteq[K]^3$ and $c$ is constantly $1$ on $[K]^3$, it is also constantly $1$ on $[K']^3$.

Thus $\kappa\to(\beta',n')_2^3$ holds.

(b) Fix a coloring $c':[\kappa']^3\to\{0,1\}$. Restrict $c'$ to the subset $[\kappa]^3\subseteq[\kappa']^3$ (identifying $\kappa$ with the initial segment $\{\alpha<\kappa\}\subseteq\kappa'$). Applying $\kappa\to(\beta,n)_2^3$ to this restricted coloring yields either a $0$-homogeneous $H\subseteq\kappa$ of order type $\beta$, or a $1$-homogeneous $K\subseteq\kappa$ of size $n$. In either case, viewing $H$ or $K$ as subsets of $\kappa'$ gives the required homogeneous set for $c'$. Hence $\kappa'\to(\beta,n)_2^3$.
\hfill$\square$

\medskip
\noindent\textbf{Corollary 70.1 (what the recorded Erd\H{o}s--Rado theorem immediately implies).}
For each integer $m\ge 2$, the recorded statement $\mathfrak c\to(\omega+m,4)_2^3$ implies
\[\mathfrak c\to(\beta,4)_2^3\quad\text{for every ordinal }\beta\le \omega+m.
\]

\noindent\textbf{Proof.}
Apply Lemma 70.2(a) with $\beta'\le\omega+m$ and $n'=4$.\hfill$\square$

\medskip
\noindent\textbf{FAST REALITY CHECK (small parameter sanity).}
For this partition relation, the only genuinely nontrivial finite parameter is $n\ge 4$. Lemma 70.1 shows the statement is automatic for $n=2$ and reduces to a trivial dichotomy for $n=3$. Thus, in the main question, the first new difficulty begins at $n=4$ and countable $\beta$ beyond $\omega+\text{(finite)}$.

\medskip
\noindent\textbf{VERIFICATION.}
\begin{itemize}
\item Lemma 70.1(a): checked that $[K]^3$ is empty for $|K|=2$, so ``$1$-homogeneous'' is vacuous.
\item Lemma 70.1(b): checked the dichotomy ``some triple is colored $1$'' vs ``all triples colored $0$''; in the latter case, the existence of a subset of order type $\beta$ uses only $\beta\le\kappa$.
\item Lemma 70.2: verified both parameter reductions by subset selection and restriction to an initial segment.
\end{itemize}

\medskip
\noindent\textbf{FINAL: \textbf{UNRESOLVED}.}
\begin{itemize}
\item[(i)] \emph{Strongest proved partial result.} The relation holds for all countable $\beta$ when $n\in\{2,3\}$ (Lemma 70.1). Moreover, the recorded Erd\H{o}s--Rado theorem implies $\mathfrak c\to(\beta,4)_2^3$ for every $\beta\le \omega+m$ (for each fixed $m\ge 2$) by monotonicity (Corollary 70.1).
\item[(ii)] \emph{First gap (crisp).} Decide whether $\mathfrak c\to(\omega^2,4)_2^3$ holds (already beyond the $\omega+\text{finite}$ range stated in the problem).
\item[(iii)] \emph{Top 3 next moves.}
(1) Attempt an induction on countable ordinals $\beta$ proving $\mathfrak c\to(\beta,4)_2^3$, identifying exactly where the Erd\H{o}s--Rado proof for $\omega+n$ breaks.
(2) Try to construct an explicit $2$-coloring of $[\mathfrak c]^3$ that avoids $0$-homogeneous sets of some fixed order type (e.g. $\omega^2$) while also avoiding $1$-homogeneous $4$-sets.
(3) Explore whether additional axioms (e.g. CH) affect the statement by testing known ``canonical'' colorings on $[\omega_1]^3$ and attempting to lift them to size $\mathfrak c$.
\item[(iv)] \emph{Minimal counterexample structure.} A counterexample would be a coloring $c:[\mathfrak c]^3\to\{0,1\}$ such that (a) every subset of $\mathfrak c$ of order type $\beta$ contains a triple colored $1$ (so it is not $0$-homogeneous), and simultaneously (b) every $n$-element subset contains a triple colored $0$ (so it is not $1$-homogeneous).
\end{itemize}


% END REUSED SOURCE: attacks/erdos/gpt_pro_5.2/70.tex

\subsection*{COMPLETION ESTIMATE}
The primary source supplies no extractable numerical completion estimate. Its mathematical progress is not numerically assessed here.

Newly verified progress in this collection pass: $0\%$. This measures new verification work, not the amount of mathematics achieved in the reused text.
